\documentclass[10pt]{ctexart}
\usepackage[a4paper,margin=1in]{geometry}
\usepackage{amsmath,amssymb,booktabs}
\usepackage{enumitem}
\usepackage{microtype}
\usepackage{hyperref}
\setlist[itemize]{leftmargin=1.5em}
\newcommand{\Whiten}{\operatorname{Whiten}}
\title{011：StepPPO-v2 改进方案}
\author{}
\date{}
\begin{document}
\maketitle
\vspace{-1.2em}

\section*{目标}
本文记录 StepPPO-v2 的第一轮改进。改进目标不是重新定义 StepPPO，而是在当前 v1 的不稳定性和运行时间诊断基础上，做两个低风险修正：降低 value loss 对 actor backbone 的干扰，并减少 old-log-probability 阶段的重复 forward 开销。

\section*{问题回顾}
StepPPO-v1 的主要问题包括：
\begin{itemize}
\item valid action ratio 明显低于 GiGPO，说明 action format 退化。
\item response clipping 更高，导致 rollout generation 和后续 forward 更慢。
\item value loss 较大且波动，且默认会回传到 shared actor backbone。
\item final advantage 由两个 whitened signals 平均得到，可能削弱 policy-gradient scale。
\item old-log-probability 阶段额外计算 state value，带来明显运行开销。
\end{itemize}

\section*{改进 1：Detach Value Backbone}
原 v1 配置为
\[
\texttt{detach\_value\_backbone=False},
\qquad
\texttt{value\_loss\_coef=0.1}.
\]
这会让 noisy value regression 直接更新 actor backbone。v2 改为
\[
\texttt{detach\_value\_backbone=True},
\qquad
\texttt{value\_loss\_coef=0.03}.
\]
这样 value loss 主要训练 value head，而不再直接牵引语言模型 backbone。policy loss 仍然正常更新 actor。该改动的目标是减少 invalid action、response clipping 和 gradient norm 震荡。

\section*{改进 2：Final Advantage Whitening}
v1 的 advantage 为
\[
A_t^{\mathrm{final,v1}}
=\frac{\widehat A_t^{\mathrm{epi}}+w\widehat A_t^{\mathrm{step}}}{1+w}.
\]
由于 $w=1$ 时相当于两个 normalized signals 的平均，final advantage 的尺度可能偏小。v2 保留该混合形式，但在混合后增加可配置 whitening：
\[
\bar A_t
=\frac{\widehat A_t^{\mathrm{epi}}+w\widehat A_t^{\mathrm{step}}}{1+w},
\qquad
A_t^{\mathrm{final,v2}}=\Whiten(\bar A_t).
\]
对应配置为
\[
\texttt{algorithm.step\_ppo.normalize\_final\_advantage=True}.
\]
默认配置中该开关保持关闭，以保证 v1/v0 的行为不变。

\section*{改进 3：Old Log Probability 与 Value 单 Forward}
原实现中，当 value head 启用时，old-log-probability 阶段先调用普通 actor forward 计算 token log probability，再额外调用一次 value forward 计算 old state values。v2 将该路径改为调用已有的 combined forward：
\[
(\log\pi_{\theta_{\mathrm{old}}}, V^{\mathrm{old}}_t)
= f_{\theta_{\mathrm{old}},\psi_{\mathrm{old}}}(s_t,a_t).
\]
这不改变算法输出，只减少重复计算。预期主要降低 \path{timing_s/old_log_prob}。

\section*{改进 4：Episode Advantage 与 GiGPO 对齐}
GiGPO 的 episode advantage 不是对整个 batch 做全局 whitening，而是在同一个 initial task group 内做 relative normalization。设一个训练 batch 中共有 $N$ 个 action/step samples，第 $i$ 个 sample 的符号定义如下：
\begin{itemize}
\item $u_i$：episode group id，对应代码中的 \path{uid}。同一个 $u_i$ 表示同一个 initial task/prompt 的多次 rollout。
\item $\tau_i$：trajectory id，对应代码中的 \path{traj\_uid}。同一个 $\tau_i$ 表示同一条 sampled episode 中的不同 step samples。
\item $\ell\in\{1,\ldots,L\}$：当前 action response 的 token 位置。
\item $m_{i,\ell}\in\{0,1\}$：response mask，只在有效 action tokens 上参与 policy loss。
\item $r^{\mathrm{tok}}_{i,\ell}$：第 $i$ 个 sample 在 token $\ell$ 上的 reward 张量值。
\item $S_i$：第 $i$ 个 sample 对应的 episode-level scalar score，定义为
\[
S_i=\sum_{\ell=1}^{L} r^{\mathrm{tok}}_{i,\ell}.
\]
\end{itemize}

对每个 group $u$，定义当前 GiGPO 代码默认使用的 row-level group set：
\[
\mathcal I_u=\{i\in\{1,\ldots,N\}:u_i=u\},
\qquad
n_u=|\mathcal I_u|.
\]
当 \path{compute\_mean\_std\_cross\_steps=True} 时，$\mathcal I_u$ 包含该 group 下所有 step rows。因此如果一条 trajectory 更长，它会在 episode advantage 的 group statistics 中贡献更多 rows。这是当前 GiGPO 实现的默认行为。

若 $n_u>1$，定义 group mean 和 group standard deviation 为
\[
\mu_u=\frac{1}{n_u}\sum_{j\in\mathcal I_u}S_j,
\]
\[
\sigma_u=\sqrt{\frac{1}{n_u-1}\sum_{j\in\mathcal I_u}(S_j-\mu_u)^2}.
\]
这里 $\sigma_u$ 对齐当前实现中 \path{torch.std} 的默认 sample standard deviation。若 $n_u=1$，当前 GiGPO episode 实现采用
\[
\mu_u=0,
\qquad
\sigma_u=1.
\]

GiGPO 支持两种 episode normalization mode。对于 \path{mean\_norm}，只做 group mean centering：
\[
A_i^{\mathrm{epi}}
=
S_i-\mu_{u_i}.
\]
对于 \path{mean\_std\_norm}，同时除以 group standard deviation：
\[
A_i^{\mathrm{epi}}
=
\frac{S_i-\mu_{u_i}}{\sigma_{u_i}+\epsilon}.
\]
最后将该 scalar episode advantage broadcast 到当前 action 的有效 response tokens：
\[
A_{i,\ell}^{\mathrm{epi}}
=
A_i^{\mathrm{epi}}\,m_{i,\ell}.
\]

为了让 StepPPO 的 episode term 与 GiGPO 对齐，v2 不应再对 $A_i^{\mathrm{epi}}$ 做额外的 batch-level whitening：
\[
A_i^{\mathrm{epi,StepPPO-v2}}
=
A_i^{\mathrm{epi,GiGPO}},
\qquad
\text{no extra }\Whiten_{\mathcal B}(\cdot).
\]
WebShop 的 GiGPO 脚本默认使用 \path{mode=mean\_norm}，因此对应的对齐版本应默认采用
\[
A_i^{\mathrm{epi,align}}
=
S_i-\mu_{u_i}.
\]

需要特别区分的是，一个更严格的 trajectory-level episode definition 可以先对同一 $(u,\tau)$ 去重：
\[
\widetilde{\mathcal I}_u
=
\left\{\min\{i:u_i=u,\tau_i=\tau\}:\tau\in\mathcal T_u\right\},
\]
再用 $\widetilde{\mathcal I}_u$ 计算 $\mu_u$ 和 $\sigma_u$。这个版本在语义上更接近“每条 episode 只贡献一次”，但它不等价于当前 GiGPO 默认代码行为。因此若实验目标是严格复现 GiGPO baseline，应优先采用 row-level group set $\mathcal I_u$；若目标是更干净的 StepPPO 变体，可以把 trajectory-level 去重作为单独 ablation。



\section*{新增脚本}
正式 StepPPO-v2 脚本为 \path{exps/run_webshop_step_ppo_v2_4gpu_paper_align_simple.sh}。旧的 \path{v2\_stable} 脚本已删除，后续只保留当前确定版 v2 入口。
核心覆盖配置为：
\begin{verbatim}
algorithm.step_ppo.step_advantage_w=0.5
algorithm.step_ppo.episode_mode=mean_norm
algorithm.step_ppo.final_advantage_mode=direct
algorithm.step_ppo.normalize_episode_advantage=False
algorithm.step_ppo.normalize_step_advantage=True
algorithm.step_ppo.normalize_final_advantage=False
actor_rollout_ref.actor.value_head.detach_value_backbone=False
actor_rollout_ref.actor.value_head.value_loss_coef=0.03
actor_rollout_ref.actor.ppo_micro_batch_size_per_gpu=8
trainer.experiment_name=step_ppo_v2_qwen2.5_1.5b_4gpu_paper_align_full_e250
\end{verbatim}

\section*{建议运行命令}
单 seed 首先运行：
\begin{verbatim}
bash exps/run_webshop_step_ppo_v2_4gpu_paper_align_simple.sh env.seed=2026
\end{verbatim}
如需 debug，可以额外覆盖 \path{trainer.total_epochs} 或 \path{trainer.test_freq}，但正式 wall-clock 对比不应改变 validation protocol。

\section*{需要观察的指标}
v2 是否有效，优先看以下指标：
\begin{itemize}
\item \path{episode/valid_action_ratio} 是否接近 GiGPO。
\item \path{response_length/clip_ratio} 是否下降。
\item \path{actor/grad_norm} 是否低于 v1 并在后期下降。
\item \path{actor/value_loss} 是否下降且不压制 policy learning。
\item \path{timing_s/old_log_prob} 是否低于早期 v1 实现。
\item \path{timing_s/update_actor} 是否因 micro-batch 对齐 GiGPO 而下降。
\item validation task score、success rate、text score 是否优于 v1。
\end{itemize}

\section*{结论}
StepPPO-v2 是 GiGPO-aligned 的 StepPPO 版本。它保留 StepPPO 的核心定义：step-wise GAE、shared actor value head、episode-level relative correction；同时将 episode advantage 与 GiGPO 的 \path{mean\_norm} 对齐，将 final advantage 改为 direct addition，降低 value loss coefficient，并将 actor PPO micro-batch 对齐 GiGPO。value loss 仍更新 shared actor backbone，即 \path{detach\_value\_backbone=False}。

\end{document}
